\documentclass[twoside, 11pt, a4paper]{article}

% --- 基础宏包 ---
% \usepackage[utf8]{inputenc}         % 支持UTF-8编码 (现代TeX引擎默认支持，通常不需要)
\usepackage[T1]{fontenc}            % 改进字体渲染
\usepackage[UTF8]{ctex}             
% Windows系统下使用UTF-8编码的中文支持，提供中文字体和排版功能（默认中文字为宋体）
% macOS系统下使用UTF-8编码的中文支持，提供中文字体和排版功能（默认中文字为华文宋体）
% Linux系统下使用UTF-8编码的中文支持，提供中文字体和排版功能（默认中文字为文泉驿正黑）
\usepackage{tocloft}                % 目录格式控制  
\usepackage{titlesec}               % 标题格式控制
\usepackage{geometry}               % 页面布局控制
\usepackage{fontspec}               % 允许调用系统字体
\usepackage{xcolor}                 % 颜色支持
\usepackage{amsmath, amssymb}       % TA必备：处理矩阵、向量、渲染方程
\usepackage{bm}                     % 加粗数学符号（常用于表示向量）
\usepackage{listings}               % 核心代码显示包
\usepackage{indentfirst}            % 强制首行缩进
\usepackage{graphicx}               % 插入渲染截图
\usepackage{hyperref}               % 启用文档内及外部链接跳转
\usepackage{zhnumber}               % 支持中文数字（如“第一章”中的“第一”）
\usepackage{caption}
\usepackage{bicaption}              % 核心：支持中英文双语题注
\usepackage{subcaption}             % 支持分图 (a)(b)
\usepackage{booktabs}               % 核心：提供 \toprule, \midrule, \bottomrule 实现三线表
\usepackage{fancyhdr}               % 页眉页脚设置
\usepackage{emptypage}              % 使空白页没有页眉页脚（适用于双面打印时的空白页）
\usepackage{titling}                % 允许自定义标题页格式

% --- 页面布局 ---
\geometry{left=2.5cm, right=2.5cm, top=2.5cm, bottom=2.5cm}
\setlength{\parindent}{2em} % 设置首行缩进为2个中文字符宽度（即4个字符空格）

% --- 字号定义 (对应中国标准字号) ---
% \zihao{3} 是三号, \zihao{-3} 是小三, 以此类推
\setmainfont{Times New Roman}   % 设置西文字体为 Times New Roman
\setCJKmainfont{SimSun}[AutoFakeBold = 2.5] % 设置中文字体为宋体，并启用伪粗体（加粗程度可调）
% 字号定义
\newcommand{\erhao}{\zihao{2}}
\newcommand{\sanhao}{\zihao{3}}
\newcommand{\xiaosan}{\zihao{-3}}
\newcommand{\sihao}{\zihao{4}}
\newcommand{\xiaosi}{\zihao{-4}}
\newcommand{\wuhao}{\zihao{5}}
\newcommand{\headerfont}{\wuhao\kaishu} % 页眉字体：五号楷体

% --- 配置页眉页脚 ---
\pagestyle{fancy}
\fancyhf{} % 清除所有默认设置

% 单号页 (Odd) 居中 (Center)：显示固定文本
\fancyhead[OC]{\headerfont 硕士学位论文}

% 双号页 (Even) 居中 (Center)：显示当前一级标题
% \leftmark 通常存储最近的一个 \section 内容
% 优化页眉标题抓取，显示第x章
\renewcommand{\sectionmark}[1]{%
  \markboth{第\thesection 章\quad #1}{}%
}
\fancyhead[EC]{\headerfont \leftmark}

% 设置页码（可选，通常居中显示在页脚）
\fancyfoot[C]{\headerfont \thepage}

% 去掉页眉下方的装饰横线（如果需要留着，设为 0.4pt）
\renewcommand{\headrulewidth}{0.4pt}

% 优化标题提取样式
% 默认的 \leftmark 带有“第x节”编号，如果你只想显示标题文字，取消下一行注释
% \renewcommand{\sectionmark}[1]{\markboth{#1}{}}

% --- 目录样式设置 ---
% 目录标题格式
\renewcommand{\contentsname}{
    \begin{center}
        \xiaosan\songti 目\qquad 录
    \end{center}
}
\renewcommand{\cftaftertoctitle}{\vspace{2em}} % 目录标题后空两行
% 目录只显示到三级标题
\setcounter{tocdepth}{3}
% 编号深度
\setcounter{secnumdepth}{3}

% 编号格式
\renewcommand{\thesection}{\zhnum{section}}
\renewcommand{\thesubsection}{\arabic{section}.\arabic{subsection}}
\renewcommand{\thesubsubsection}{\thesubsection.\arabic{subsubsection}}

% 一级标题（章）在目录中的样式：四号黑体，不缩进
\renewcommand{\cftsecfont}{\sihao\heiti}         % 目录条目文字字体
\renewcommand{\cftsecpagefont}{\sihao\heiti}     % 目录页码字体
\setlength{\cftsecindent}{0em}                   % 左侧缩进
\setlength{\cftsecnumwidth}{5em}                 % 编号宽度（留够“第1章”的空间）
\renewcommand{\cftsecleader}{\cftdotfill{\cftdotsep}} % 确保一级标题也有引导点
% 修正目录中一级标题的编号显示为“第X章”
\renewcommand{\cftsecpresnum}{第}      % 前缀只写“第”
\renewcommand{\cftsecaftersnum}{章\quad} % 后缀写“章”加空格

% 二级标题（节）在目录中的样式：小四号宋体，左空2字符
\renewcommand{\cftsubsecfont}{\xiaosi\songti}
\renewcommand{\cftsubsecpagefont}{\xiaosi\songti}
\setlength{\cftsubsecindent}{2em}                % 缩进2字符
\setlength{\cftsubsecnumwidth}{2em}

% 三级标题（小节）在目录中的样式：小四号宋体，左空4字符
\renewcommand{\cftsubsubsecfont}{\xiaosi\songti}
\renewcommand{\cftsubsubsecpagefont}{\xiaosi\songti}
\setlength{\cftsubsubsecindent}{4em}             % 缩进4字符
\setlength{\cftsubsubsecnumwidth}{3em}

% --- 论文正文格式 ---
% 一级标题（章）格式：
\titleformat{\section}
  {\centering\sanhao\songti} % 居中 (\centering)、三号 (\sanhao)、宋体 (\songti)
  {第\thesection 章}    % 第X章
  {24pt} % 编号与标题之间的距离，小四号（12pt）字号下两个字符空格的长度为24pt
  {}

% 二级标题（节）格式：
\titleformat{\subsection}
  {\sihao\songti\bfseries} % 四号 (\sihao)、宋体 (\songti)、加粗居左
  {\thesubsection }    % 1.1
  {24pt} % 编号与标题之间的距离，小四号（12pt）字号下两个字符空格的长度为24pt
  {}

% 三级标题（小节）格式：
\titleformat{\subsubsection}
  {\xiaosi\songti} % 小四 (\xiaosi)、宋体 (\songti)、居左
  {\thesubsubsection}   % 1.1.1
  {2em} % 编号与标题之间的距离，小四号（12pt）字号下两个字符空格的长度为24pt
  {}

% --- 图号按章编号，格式为 图1-1 ---
\renewcommand{\thefigure}{\arabic{section}-\arabic{figure}} 
\captionsetup{labelsep=quad} % 编号与文字间空一个全角空格

% % 定义中英文图题字体
% 中文：五号宋体 (\wuhao\songti)
% 英文：五号 Times New Roman (\wuhao\rmfamily)
\DeclareCaptionFont{chinese}{\wuhao\songti}
\DeclareCaptionFont{english}{\wuhao\rmfamily}
\captionsetup[bi-first]{
    font=chinese,
    textfont=normalfont,
    name={图},
}
\captionsetup[bi-second]{
    font=english,
    textfont=normalfont,
    name={Fig.},
}

% --- 表号按章编号，格式为 表1-1 ---
\renewcommand{\thetable}{\arabic{section}-\arabic{table}} 

% 定义表格题注样式
\DeclareCaptionFont{chinese}{\wuhao\songti}
\DeclareCaptionFont{english}{\wuhao\rmfamily}
\captionsetup[table][bi-first]{
    font=chinese,
    name={表},
    labelsep=quad
}
\captionsetup[table][bi-second]{
    font=english,
    name={Table},
    labelsep=quad
}
% 设置题注在表格上方（符合国家标准）
\captionsetup[table]{position=top, skip=5pt}

% --- 设计与色彩 ---
\definecolor{codeblue}{RGB}{45, 156, 219}
\definecolor{codegreen}{RGB}{34, 139, 34}
\definecolor{backcolour}{RGB}{245, 245, 245}

% --- 代码块配置 (listings) ---
\lstset{
    backgroundcolor=\color{backcolour},  % 背景色
    commentstyle=\color{codegreen},      % 注释颜色
    keywordstyle=\color{codeblue},       % 关键字颜色
    numberstyle=\tiny\color{gray},       % 行号样式
    basicstyle=\ttfamily\small,          % 等宽字体
    breaklines=true,                     % 自动换行
    captionpos=b,                        % 标题位置：底部
    frame=single,                        % 代码框样式
    showstringspaces=false               % 不显示字符串中的空格
    extendedchars=false,                 % 解决代码跨页可能导致的编译错误
    keepspaces=true,                     % 保持空格
    columns=flexible,                    % 解决代码中英混排时间距过大的问题
}

% --- 图片与超链接 ---
\hypersetup{colorlinks=true, linkcolor=black, urlcolor=cyan}

% --- 文档信息 ---
% 变量赋值，不会显示在文档中，但可以在后续使用
\title{\textbf{技术美术（TA）学习笔记}} % textbf{} 加粗
\author{Technical Artist}
\date{2026年3月22日}

% -------------------- 正文开始 --------------------
\begin{document}
\xiaosi % 这里开启全局小四号字体设置，后续正文内容将默认使用小四号字体，除非在局部重新定义。
% -------------------- 标题页开始 
\begin{titlepage}
    \pagenumbering{gobble} % 这一页不计页码
    \vspace*{\fill}       % 顶部弹簧
    
    \begin{center}
        {\sanhao\bfseries 技术美术（TA）学习笔记 \par} % 标题
        \vspace{2em}
        {\sihao Technical Artist \par}               % 作者
        \vspace{1em}
        {\sihao 2026年3月22日 \par}                  % 日期
    \end{center}
    
    \vspace*{\fill}       % 底部弹簧
\end{titlepage}
% -------------------- 标题页结束



% -------------------- 中文摘要部分开始 
\cleardoublepage % 确保摘要页从右侧（单数页）开始
\pagenumbering{Roman} % 页码：大写罗马 I, II...
\phantomsection % 创建一个锚点，使得目录中的链接能够正确跳转到摘要页
\addcontentsline{toc}{section}{摘要} % 将“摘要”加入目录，作为一级标题（section）显示
% \begin{center}
% \end{center}
% 将其内部包含的所有内容（文字、图片、公式等）在页面水平方向上居中对齐。
\begin{center}
    % “摘要”二字（小三号宋体），摘要二字间用二个字符空格分开
    {\xiaosan\songti 摘\qquad 要} % 小三号宋体，中间两个全角空格（\quad 是一个字符宽度、\qquad 是两个字符宽度）
\end{center}

\vspace{1em} % 下空一行

{ % ******摘要内容（默认小四号宋体）
本文档内容主要涵盖了二次元（NPR）渲染核心、引擎底层与性能优化、数字内容创作（DCC）工具与自动化、数学与逻辑基础等方面的知识点，旨在为技术美术（TA）提供一个系统的学习笔记。
}

\vspace{1em} % 下空一行

%左对齐打印“关键词”三字（五号宋体加粗）
\noindent % 关键词行取消缩进，实现左对齐
{\wuhao\textbf {关键词：}} % 五号宋体加粗
{\wuhao\songti 渲染管线\qquad 材质Shader\qquad 性能优化\qquad HLSL算法} % 关键词（五号宋体）每一关键词后之间用二个字符空格分开，最后一个关键词后不打标点符号

% -------------------- 中文摘要部分结束 



% -------------------- 英文摘要部分开始 
\cleardoublepage % 确保英文摘要页从右侧（单数页）开始
\pagenumbering{roman} % 页码：小写罗马 i, ii...
\phantomsection % 创建一个锚点，使得目录中的链接能够正确跳转到英文摘要页
% \addcontentsline{toc}{section}{ABSTRACT} % 将“ABSTRACT”加入目录，作为一级标题（section）显示

\begin{center}
    % ABSTRACT (大写), Times New Roman 四号加粗
    {\sihao\bfseries ABSTRACT} 
\end{center}

\addcontentsline{toc}{section}{\protect\rmfamily\bfseries ABSTRACT}


\vspace{1em} % 下空一行

{% ******英文摘要内容：首行缩进四个字符空格，Times New Roman 小四号字体
\xiaosi The content of this document mainly covers knowledge points in aspects such as the core of Non-Photorealistic Rendering(NPR), the engine bottom layer and performance optimization, digital content creation (DCC) tools and automation, and mathematical and logical foundations. It aims to provide a systematic learning notebook for technical artists (TA).
}

\vspace{1em} % 下空一行

% 3. KEYWORDS：小四号加粗，关键词内容不加粗
\noindent 
{\xiaosi\bfseries KEYWORDS:} 
{\xiaosi Rendering Pipeline；Shader Material；Performance Optimization；HLSL Algorithm}
% 每个关键词首字母大写，用全角分号“；”隔开，末尾无标点
% -------------------- 英文摘要部分结束 



% -------------------- 目录页开始 --------------------
\cleardoublepage
\pagestyle{fancy}   
\pagenumbering{gobble}

% 1. 手动强制更新标记，确保 \leftmark 内容为“目  录”
\markboth{目\qquad 录}{目\qquad 录}

% 2. 针对目录页，将偶数页居中也改为“目  录”
% 这样无论目录有几页，双号页都会显示“目  录”
\fancyhead[EC]{\headerfont 目\quad 录} 

\addtocontents{toc}{\protect\thispagestyle{fancy}}
\tableofcontents

\cleardoublepage
% 3. 目录结束，切回正文前，把偶数页页眉改回动态抓取模式
\fancyhead[EC]{\headerfont \leftmark}

% -------------------- 目录页结束 --------------------



% -------------------- 正文内容部分开始 

\pagenumbering{arabic} % 页码：阿拉伯数字 1, 2...

% ******第一章
\section{二次元核心渲染}

% ******第一节
\subsection{核心渲染算法}

% % 图片示例：插入一张名为“1-1.png”的图片，宽度占页面宽度的80%，并添加双语图题
% \begin{figure}[htbp]
%     \centering
%     \includegraphics[width=0.8\textwidth]{1-1.png} % 图片文件名为“1-1.png”，请确保该图片文件与.tex文件在同一目录下，或者提供正确的路径
    
%     % 双语图题：{中文}{英文}
%     \bicaption{渲染管线流程图}{Flowchart of rendering pipeline}
    
%     % 在图题下方标注文献来源，如有取消注释
%     % \centering{\small 来源：参考文献[12]} 
% \end{figure}

% % 表格示例：
% \begin{table}[htbp]
%     \centering
%     % 中英文对照标题
%     \bicaption{Unity 常用渲染路径性能对比}{Performance comparison of common rendering paths in Unity}
%     \label{tab:render_compare}
    
%     \begin{tabular}{@{} lcc @{}} % @{} 用于去除表格两端的额外边距
%         \toprule
%         渲染路径 (Path) & 动态光源支持 & 内存开销 (VRAM) \\
%         \midrule
%         前向渲染 (Forward) & 较少 & 低 \\
%         延迟渲染 (Deferred) & 极多 & 高 \\
%         分块渲染 (Forward+) & 较多 & 中 \\
%         \bottomrule
%     \end{tabular}
% \end{table}

% ******第一小节
\subsubsection{多分色阴影}

{% ******文章内容：

}


\end{document}